\section{Ukratko}

\begin{kljucno}[Osam poruka za brzo čitanje]
\begin{enumerate}[leftmargin=1.6em]
\item \textbf{Što je $\alpha$.} $\alpha$ je eksponent kojim efektivni kapacitet izmjene između transportnog kompartimenta i tkiva raste sa \emph{strukturnim} volumenom, $A\propto V^{\alpha}$. Kapacitet izmjene je površina × gustoća transportera × propusnost (H2). $\alpha$ \emph{nije} metabolički eksponent $b$ (pogl.~\ref{sec:definicija}, \ref{sec:zamke}).
\item \textbf{Raspon.} Raspon $2/3\le\alpha\le1$ dobar je radni prior za odrasle organizme s unutarnjom mrežom, ali nije zakon. Vrijednosti ispod $2/3$ zabilježene su, a iznad $1$ javljaju se prolazno, dok se mreža gradi (T11, E3, E4).
\item \textbf{MTE ne daje $\alpha=1$.} Prema West–Brown–Enquistu (WBE) $M^{3/4}$ je \emph{gornja} granica efektivne površine izmjene fraktalne mreže koja ispunjava volumen. Prema 1 se eksponent približava tek uz čisto Poiseuilleovo grananje (T8, T9). Ni fraktalna dimenzija izmjerena na plošnim slikama mreže ne daje $\alpha=1$ (H4, P6).
\item \textbf{Morfometrija daje prior, ne procjenu.} Izmjerene površine izmjene kod sisavaca skaliraju s eksponentom $0{,}87$–$0{,}99$, iznad WBE granice. Prate vršnu potražnju za kisikom, a ne bazalni metabolizam (E1, E2).
\item \textbf{Rast nikad nije dovoljan.} Uz brz transportni kompartiment somatski tok je $\kappa\,\pd_A$ s $\kappa=\beta_S/(\beta_S+\beta_R)$, a $L_\infty$ ne ovisi o $\alpha$ (P1). U Monte Carlu procjene iz same krivulje rasta raspršene su po cijelom dopuštenom rasponu (za istinski $\alpha=0{,}75$ 10.–90. percentil je \MCAsevenfiveLo–\MCAsevenfiveHi, a RMSE do \MCArmse; S2, S6).
\item \textbf{Gladovanje je najinformativnije, spremište drugo.} Specifična brzina pražnjenja spremišta pri gladovanju skalira kao $V^{\alpha-1}$ i ne ovisi o eksponentu asimilacije $a$ (P3). Serija gladovanja u tri veličine daje RMSE do \MCErmse, najbolje od svih dizajna. Gustoća spremišta po strukturi skalira kao $V^{a-\alpha}$ (P2) i uz 8–16 mjerenja daje RMSE do \MCBrmse–\MCCrmse{}, ovisno o mjernoj pogrešci (S6).
\item \textbf{Pogrešan $a$ spušta $\hat\alpha$.} Ako je stvarni $a=0{,}72$, a fit pretpostavlja $2/3$, medijan $\hat\alpha$ pada s \MCCokMed{} na \MCCmisMed{} (samo spremište) i s \MCFokMed{} na \MCFmisMed{} (spremište i gladovanje), u skladu s $\hat\alpha\approx\alpha-(a-2/3)$ iz P2 (S6). Pristranost ostaje i uz gladovanje, pa unos hrane prema veličini treba mjeriti.
\item \textbf{Protokol:} prior iz tipa mreže → gladovanje prema veličini, spremište po strukturi prema veličini i unos hrane prema veličini → zajednički fit dinamičkog modela s profilom vjerodostojnosti → filogenetski test (H6; pogl.~\ref{sec:protokol}). Nagib spremišta ne čitati „na oko“: naivna formula podcjenjuje $\alpha$ za \NaiveBiasMin–\NaiveBiasMax{} (S4).
\end{enumerate}
\end{kljucno}
